% answers_short.tex —— 1986 年·上海卷·文史类 · 简短答案速查
% 由 tools/make_short_answers.py 自 answers.tex 的 \gksolhead/\gkanswer 提取，
% 并与 1986/1986_数学真题_上海文_答案要点.md 逐题比对后生成，勿手改（改源 MD 后重跑）。

\begin{gknotice}
\textbf{使用说明}：本节只给最终答案便于快速核对，完整推导见解析版。答案为编写者独立推导，\textbf{非官方标准答案}。
\end{gknotice}

\gkpart{大题}
\begin{itemize}[leftmargin=2.6em]
  \item[\textbf{1.}] (1) B；(2) B；(3) D；(4) B；(5) C；(6) A；(7) C；(8) A；(9) D；(10) B；(11) B；(12) A；(13) D；(14) B；(15) D；(16) C；(17) C；(18) D；(19) A；(20) C
  \item[\textbf{2.}] (1) $(-3,-12)$；(2) $\dfrac\pi6$；(3) $2$；(4) $1$；(5) $28$（或 $C_8^2$）；(6) $1\le x<10$；(7) $\pi$；(8) $12$；(9) $7$；(10) $(0,1)$
  \item[\textbf{3.}] $-1<x\le0$
  \item[\textbf{4.}] $\theta=\dfrac{3\pi}4$ 时最大值为 $3$；$\theta=\dfrac\pi4$ 时最小值为 $\dfrac13$
  \item[\textbf{5.}] 见解析正文。
  \item[\textbf{6.}] $90^\circ$
  \item[\textbf{7.}] $\arccos\dfrac{\sqrt3}{3}$
\end{itemize}
